% !TeX root = ../../Book3.tex
\section{Eisenstein 级数}
\label{b3:sec:2302}

前章的格不变量随基变换产生确定的幂次因子。绝对收敛是把这种格对称性转成函数变换律的条件；权二处会失去它。

\subsection{格和}
\label{b3:subsec:230201}

对整数 \(k\ge2\)，定义\term[eisenstein-jishu]{Eisenstein 级数}[Eisenstein series]
\[
G_{2k}(\tau)\coloneqq\sum_{(m,n)\in\mathbb Z^2\setminus\{(0,0)\}}(m\tau+n)^{-2k}.
\]
这就是格 \(\Lambda_\tau\) 的 \(2k\) 次倒数和。它使用全部非零格向量；若只对互素 \((m,n)\) 求和，结果须乘以 \(\sum_{r\ge1}r^{-2k}\)，不能混用两种归一化。

\subsection{收敛}
\label{b3:subsec:230202}

在任意紧集 \(K\subset\mathfrak H\) 上，实线性变换 \((m,n)\mapsto m\tau+n\) 的最小奇异值有共同正下界，所以 \(|m\tau+n|\ge c_K\sqrt{m^2+n^2}\)。前章格壳估计给正规绝对收敛，故 \(G_{2k}\) 全纯。求导后的项含因子 \(m\)，但总衰减仍至少为 \((m^2+n^2)^{-k}\)，同一紧集估计允许任意固定次求导。

\ProseParagraphBreak

若 \(k=1\)，壳估计只给调和级数 \(\sum j^{-1}\)，绝对收敛失败。规定先对 \(n\) 再对 \(m\) 求和会得到有用函数，但它的变换律多出一项，不能把下面的重排论证照搬过去。

\subsection{\texorpdfstring{\(G_{2k}\)}{G2k}}
\label{b3:subsec:230203}

\begin{proposition}\label{b3:prop:23-eisenstein-transformation}
若 \(k\ge2\)、\(\gamma=\left(\begin{smallmatrix}a&b\\c&d\end{smallmatrix}\right)\in\mathrm{SL}_2(\mathbb Z)\)，则
\(G_{2k}(\gamma\tau)=(c\tau+d)^{2k}G_{2k}(\tau)\)。
\end{proposition}
\begin{proof}
每项分母为 \((m(a\tau+b)+n(c\tau+d))/(c\tau+d)\)。变换 \((m,n)\mapsto(ma+nc,mb+nd)\) 是 \(\mathbb Z^2\) 的双射，且保持非零对。绝对收敛允许依此重排，提出公共因子即得结果。
\end{proof}
令 \(\sigma_r(N)=\sum_{d\mid N}d^r\)。从 \(\pi\cot\pi z=-\pi\mathrm{i}(1+2\sum_{r\ge1}e^{2\pi\mathrm{i}rz})\) 在上半平面的展开出发，求 \(2k-1\) 次导数得到
\[
\sum_{n\in\mathbb Z}(z+n)^{-2k}
=\frac{(2\pi\mathrm{i})^{2k}}{(2k-1)!}\sum_{r\ge1}r^{2k-1}e^{2\pi\mathrm{i}rz}.
\]
先分出 \(m=0\)，将正负 \(m\) 配对，再按 \(N=mr\) 收集系数，得
\[
G_{2k}=2\zeta(2k)+\frac{2(2\pi\mathrm{i})^{2k}}{(2k-1)!}
\sum_{N\ge1}\sigma_{2k-1}(N)q^N.
\]
这里 \(\zeta(2k)\) 只表示收敛实级数。双和的绝对收敛由 \(|q|<1\) 和多项式系数增长保证，因而合并合法。

\subsection{与 \texorpdfstring{\(\wp\)}{\AnalysisPdfWeierstrassP} 的关系}
\label{b3:subsec:230204}

前章的记号现在为 \(g_2(\Lambda_\tau)=60G_4(\tau)\)、\(g_3(\Lambda_\tau)=140G_6(\tau)\)。格缩放律已经预示它们的权分别为四、六。相比之下，\(\wp(z;\Lambda_\tau)\) 还带一个曲线上的变量 \(z\)；改变 \(\tau\) 的基时必须同时缩放 \(z\)，不能将它误当成只依赖模参数的模形式。
